\section{Setup and conventions}
\label{sec:setup}

\subsection{Kinematics and definition of the cross-section}
\label{sec:kin}

We use light-cone coordinates $v^\pm=(v^0\pm v^3)/\sqrt2$ and denote transverse vectors by boldface
letters. The virtual photon is a right-mover with four-momentum
$q^\mu=(q^+,q^-,\qq)=(q^+,-Q^2/2q^+,\qq)$ and the target is a left-mover, which in the CGC effective
theory is represented by a shock wave localised at $x^+=0$. We work in the projectile light-cone gauge
$A^+=0$. The final state consists of a quark, an antiquark and a gluon, which are on-shell and whose
kinematics is summarised in table~\ref{tab:kin}. A three-momentum is denoted as $k\equiv(k^+,\kk)$ and
$\dd^3k\equiv\dd k^+\dd^2\kk$. The light-cone energies are
\begin{equation}
E_q(k)=\frac{\kk^2+\mf^2}{2k^+}\,,\qquad E_{\bar q}(p)=\frac{\pp^2+\mf^2}{2p^+}\,,\qquad
E_g(p)=\frac{\pp^2}{2p^+}\,,\qquad E_\gamma(q)=\frac{\qq^2-Q^2}{2q^+}\,.
\end{equation}

\begin{table}[t]
\centering
\begin{tabular}{lccc}
\toprule
parton & longitudinal momentum fraction & transverse momentum & mass\\
\midrule
quark & $\zeta=k_1^+/q^+$ & $\kk_1$ & $\mf$\\
antiquark & $\eta=k_2^+/q^+=1-\zeta-\xi$ & $\kk_2$ & $\mf$\\
gluon & $\xi=k_3^+/q^+$ & $\kk_3$ & $0$\\
\bottomrule
\end{tabular}
\caption{Kinematics of the final state. The transverse coordinates conjugate to $\kk_1$, $\kk_2$ and
$\kk_3$ are denoted by $\x$, $\y$ and $\z$ in the direct amplitude and by $\bx$, $\by$ and $\bz$ in the
complex-conjugate amplitude.}
\label{tab:kin}
\end{table}

Following ref.~\cite{Iancu:2022gpw}, the trijet cross-section is obtained by counting the number of
$\qqg$ triplets with given momenta in the outgoing state of the photon,
\begin{align}
&\frac{\dd\sigma^{\gT A\to\qqg+X}}{\dd^3k_1\,\dd^3k_2\,\dd^3k_3}\;(2\pi)\,\delta\big(q^+-k_1^+-k_2^+-k_3^+\big)
\nonumber\\
&\qquad=\frac12\int_{\w,\bar\w}e^{-i\qq\cdot(\w-\bar\w)}\;
{}_{\qqg}\big\langle\gamma_T^i(Q,q^+,\bar\w)\big|\,N_q(k_1)\,N_{\bar q}(k_2)\,N_g(k_3)\,
\big|\gamma_T^i(Q,q^+,\w)\big\rangle_{\qqg}\,,
\label{eq:xsdef}
\end{align}
where $N_q$, $N_{\bar q}$ and $N_g$ are the number-density operators for bare quarks, antiquarks and
gluons, $\w$ ($\bar\w$) is the transverse coordinate of the photon in the direct (complex-conjugate)
amplitude, the factor $1/2$ averages over the two transverse polarisations $i=1,2$ of the photon, and
$\int_{\x}\equiv\int\dd^2\x$. We set $\qq=0$ in what follows, without loss of generality.

The outgoing state is obtained by evolving the incoming photon with the QED and QCD interaction
Hamiltonians up to the collision time $x^+=0$, by acting with the $S$-matrix $\hat S$ of the eikonal
scattering off the shock wave, and by evolving again up to $x^+\to\infty$. To the order of interest, the
$\qqg$ component reads~\cite{Iancu:2022gpw}
\begin{align}
|\gamma\rangle_{\qqg}=|q_2\bar q_2g_2\rangle\bigg\{&
\langle q_2\bar q_2g_2|\hat S-1|q_1\bar q_1g_1\rangle\,
\frac{\langle q_1\bar q_1g_1|H_{\rm QCD}|\qqbar\rangle}{E_{q_1\bar q_1g_1}-E_\gamma}
\nonumber\\
&+\frac{\langle q_2\bar q_2g_2|H_{\rm QCD}|q_1\bar q_1\rangle}{E_{q_1\bar q_1}-E_{q_2\bar q_2g_2}}\,
\langle q_1\bar q_1|\hat S-1|\qqbar\rangle\bigg\}\,
\frac{\langle\qqbar|H_{\rm QED}|\gamma\rangle}{E_{\qqbar}-E_\gamma}
\nonumber\\
+\;|q_2\bar q_2g_2\rangle\,&\langle q_2\bar q_2g_2|\hat S-1|q_1\bar q_1g_1\rangle\,
\frac{\langle q_1\bar q_1g_1|H_{\gamma\to\qqg}|\gamma\rangle}{E_{q_1\bar q_1g_1}-E_\gamma}\,,
\label{eq:formal}
\end{align}
where a sum over the intermediate Fock states is understood. The first line describes a gluon emission
before the collision, the second line a gluon emission after the collision, and the third line the
instantaneous contribution, in which the photon decay and the gluon emission occur at the same
light-cone time. The subtraction of unity from $\hat S$ accounts for the graphs in which the photon
crosses the shock wave before decaying. The corresponding diagrams are shown in
figures~\ref{fig:qbar} and \ref{fig:q}.

\subsection{Matrix elements with massive quarks}
\label{sec:me}

We need three matrix elements. They are written in terms of two-component helicity spinors $\chi_\lambda$,
$\lambda=\pm1/2$, and of the Pauli matrices $\sigma^i$; $\ee^{ij}$ is the antisymmetric tensor in two
dimensions with $\ee^{12}=1$, $e_f$ is the electric charge of the flavour $f$ in units of $e$, and
$\alpha,\beta$ ($a,b$) are colour indices in the fundamental (adjoint) representation.

\paragraph{Photon decay.} For a transverse photon with polarisation index $i$ decaying into a quark with
momentum $k$ and an antiquark with momentum $p$,
\begin{align}
\big\langle q^\alpha_{\lambda_1}(k)\,\bar q^\beta_{\lambda_2}(p)\big|H_{\gamma\to\qqbar}\big|\gamma_T^i(q)\big\rangle
={}&(2\pi)^3\delta^{(3)}(q-k-p)\,\frac{e\,e_f\,\delta_{\alpha\beta}}{2\sqrt{2q^+}}
\nonumber\\
&\times\chi^\dagger_{\lambda_1}\bigg[\frac{2q^i}{q^+}
-\frac{\bm\sigma\cdot\kk-\mf\sigma^3}{k^+}\,\sigma^i
-\sigma^i\,\frac{\bm\sigma\cdot\pp+\mf\sigma^3}{p^+}\bigg]\chi_{\lambda_2}\,.
\label{eq:Hgamma}
\end{align}
With $k^+=\zeta q^+$, $p^+=(1-\zeta)q^+$ and $\kt\equiv\kk-\zeta\qq$, the spinor structure reduces to
\begin{equation}
\frac{2q^i}{q^+}-\frac{\bm\sigma\cdot\kk-\mf\sigma^3}{k^+}\,\sigma^i
-\sigma^i\,\frac{\bm\sigma\cdot(\qq-\kk)+\mf\sigma^3}{q^+-k^+}
=\frac{\big[(2\zeta-1)\delta^{ij}+i\ee^{ij}\sigma^3\big]\kt^j+\mf\,\sigma^3\sigma^i}{\zeta(1-\zeta)\,q^+}\,,
\end{equation}
where we used $\sigma^i\sigma^j=\delta^{ij}+i\ee^{ij}\sigma^3$ and $\sigma^i\sigma^3=-\sigma^3\sigma^i$.
Taking the matrix element between helicity states, we obtain
\begin{align}
\big\langle q^\alpha_{\lambda_1}(k)\,\bar q^\beta_{\lambda_2}(p)\big|H_{\gamma\to\qqbar}\big|\gamma_T^i(q)\big\rangle
=(2\pi)^3\delta^{(3)}(q-k-p)\;
\frac{e\,e_f\,\delta_{\alpha\beta}\,\big[\varphi^{ij}_{\lambda_1\lambda_2}(\zeta)\,\kt^j
+\mf\,\mu^i_{\lambda_1\lambda_2}\big]}{2\sqrt{2(q^+)^3}\;\zeta(1-\zeta)}\,,
\label{eq:Hgamma2}
\end{align}
with the helicity-conserving and helicity-flip structures
\begin{align}
\varphi^{ij}_{\lambda_1\lambda_2}(\zeta)&\equiv
\chi^\dagger_{\lambda_1}\big[(2\zeta-1)\delta^{ij}+i\ee^{ij}\sigma^3\big]\chi_{\lambda_2}
=\delta_{\lambda_1\lambda_2}\big[(2\zeta-1)\delta^{ij}+2i\ee^{ij}\lambda_1\big]\,,
\label{eq:phi}\\
\mu^i_{\lambda\lambda'}&\equiv\chi^\dagger_\lambda\,\sigma^3\sigma^i\,\chi_{\lambda'}
=\delta_{\lambda',-\lambda}\big[2\lambda\,\delta^{i1}-i\,\delta^{i2}\big]\,,
\qquad \sum_{i}\mu^i_{\lambda\lambda'}\,\mu^{i\,*}_{\lambda\lambda'}=2\,\delta_{\lambda',-\lambda}\,.
\label{eq:mu}
\end{align}
The mass term flips the helicity, $\lambda_2=-\lambda_1$, and carries no transverse momentum.

\paragraph{Gluon emission by the antiquark.} For an antiquark with momentum $p$ and helicity
$\lambda_2$ which emits a gluon with momentum $q'$, polarisation $i'$ and colour $a'$,
\begin{align}
&\big\langle q^{\alpha'}_{\lambda_1'}(k')\,\bar q^{\beta'}_{\lambda_2'}(p')\,g^{a'}_{i'}(q')\big|
H_{\bar q\to\bar qg}\big|q^\alpha_{\lambda_1}(k)\,\bar q^\beta_{\lambda_2}(p)\big\rangle
=-(2\pi)^6\,\delta^{(3)}(k-k')\,\delta^{(3)}(p-p'-q')\;\delta_{\alpha'\alpha}\,\delta_{\lambda_1\lambda_1'}
\nonumber\\
&\hspace{2.5cm}\times\frac{g\,t^{a'}_{\beta\beta'}}{2\sqrt{2q'^+}}\;
\chi^\dagger_{\lambda_2}\bigg[\frac{2q'^{i'}}{q'^+}
-\sigma^{i'}\,\frac{\bm\sigma\cdot\pp'+\mf\sigma^3}{p'^+}
-\frac{\bm\sigma\cdot\pp+\mf\sigma^3}{p^+}\,\sigma^{i'}\bigg]\chi_{\lambda_2'}\,.
\label{eq:Hg}
\end{align}
For $p^+=(1-\zeta)q^+$, $p'^+=\eta q^+$, $q'^+=\xi q^+$ and with the relative transverse momentum
$\pt\equiv\qq'-\frac{\xi}{1-\zeta}\,\pp$, the terms proportional to $\pp$ cancel and
\begin{equation}
\chi^\dagger_{\lambda}\bigg[\frac{2q'^{i'}}{q'^+}
-\sigma^{i'}\,\frac{\bm\sigma\cdot\pp'+\mf\sigma^3}{p'^+}
-\frac{\bm\sigma\cdot\pp+\mf\sigma^3}{p^+}\,\sigma^{i'}\bigg]\chi_{\lambda_2}
=\frac{\tau^{i'n}_{\lambda\lambda_2}(\xi,\eta)\,\pt^{\,n}+\xi\,\of\,\mu^{i'}_{\lambda\lambda_2}}{\xi\,\eta\,q^+}\,,
\label{eq:Hg2}
\end{equation}
where
\begin{equation}
\tau^{mn}_{\lambda\lambda_2}(\xi,\eta)\equiv\delta_{\lambda\lambda_2}\Big[\big(2(1-\zeta)-\xi\big)\delta^{mn}
+2i\,\xi\,\ee^{mn}\lambda\Big]\,,\qquad
\of\equiv\frac{\xi\,\mf}{1-\zeta}\,.
\label{eq:tau}
\end{equation}
The two arguments of $\tau$ are the longitudinal fractions of the daughter gluon and antiquark, and
the coefficient of $\delta^{mn}$ is the sum of the fractions of the antiquark before and after the
emission. The mass term follows from
$\mf\,\sigma^3\sigma^{i'}\big[1/p'^+-1/p^+\big]=\xi\,\of\,\sigma^3\sigma^{i'}/(\xi\eta q^+)$. It is
suppressed by two powers of the gluon fraction with respect to $\mf$: this is the origin of the
dead-cone scale $\of$, which is nothing but the quark mass times the splitting fraction
$\xi/(1-\zeta)$ of the gluon.

\paragraph{Instantaneous vertex.} The instantaneous part of the fermion propagator is proportional to
$\gamma^+/2k^+$ and is independent of the mass. The corresponding matrix element therefore keeps its
massless form~\cite{Iancu:2022gpw},
\begin{align}
\big\langle q^\alpha_{\lambda_1}(k)\,\bar q^\beta_{\lambda_2}(s)\,g^b_j(p)\big|H_{\gamma\to\qqg}\big|\gamma_T^i(q)\big\rangle
={}&-(2\pi)^3\delta^{(3)}(q-k-p-s)
\nonumber\\
&\times\frac{g\,e\,e_f\,\delta_{\lambda_1\lambda_2}\,t^b_{\alpha\beta}}{4\sqrt{q^+p^+}}\;
\frac{\delta^{ij}+2i\ee^{ij}\lambda_1}{p^++s^+}\,,
\label{eq:Hinst}
\end{align}
for the graph in which the instantaneous propagator is that of the antiquark; for the quark, the
numerator becomes $-(\delta^{ij}-2i\ee^{ij}\lambda_1)$ and the denominator $p^++k^+$.

The two-component expressions \eqref{eq:Hgamma2}, \eqref{eq:Hg2} and \eqref{eq:Hinst}, including the
relative signs and normalisations of the mass terms and of the instantaneous terms, have been checked
against the complete amplitude computed with four-component Dirac spinors, see
appendix~\ref{app:checks}.

\subsection{Energy denominators}
\label{sec:den}

The photon decay involves the energy difference
\begin{equation}
E_{\qqbar}-E_\gamma=\frac{\kt^2+\ef^2}{2\zeta(1-\zeta)\,q^+}\,,\qquad
\ef^2\equiv\zeta(1-\zeta)\,Q^2+\mf^2\,,
\label{eq:den1}
\end{equation}
where $\ef$ is the usual mass-shifted photon scale. For the emission of the gluon by the antiquark,
\begin{align}
E_{\qqg}-E_{\qqbar}&=\frac{\pp^2}{2\xi q^+}+\frac{(\qq-\kk-\pp)^2+\mf^2}{2\eta q^+}
-\frac{(\qq-\kk)^2+\mf^2}{2(1-\zeta)q^+}
=\frac{(1-\zeta)\big(\pt^{\,2}+\of^2\big)}{2\,\xi\,\eta\,q^+}\,,
\label{eq:den2}
\end{align}
where now $\pt=\pp-\frac{\xi}{1-\zeta}(\qq-\kk)$ and $\pp$ is the gluon momentum. The mass does not
cancel in eq.~\eqref{eq:den2}, because the antiquark carries a smaller longitudinal momentum after the
emission: the gluon can never be exactly collinear to a massive emitter. The sum of
eqs.~\eqref{eq:den1} and \eqref{eq:den2} gives the denominator of the three-parton state,
\begin{equation}
E_{\qqg}-E_\gamma=\frac{\zeta(1-\zeta)^2\big(\pt^{\,2}+\of^2\big)+\xi\eta\big(\kt^2+\ef^2\big)}
{2\,\zeta(1-\zeta)\,\xi\,\eta\,q^+}\,.
\label{eq:den3}
\end{equation}
The mass-dependent constant in the numerator of eq.~\eqref{eq:den3} can be written as
\begin{equation}
\zeta(1-\zeta)^2\,\of^2+\xi\eta\,\ef^2=\zeta(1-\zeta)\,\xi\eta\,\Ef^2\,,\qquad
\Ef^2\equiv Q^2+\mf^2\left(\frac1\zeta+\frac1\eta\right)\,,
\label{eq:Ef}
\end{equation}
which shows that $\Ef^2$ plays, for the three-parton state, the role played by $Q^2$ in the massless
case. Indeed, $\Ef^2$ is the virtuality $Q^2$ shifted by the squared invariant mass
$\mf^2/\zeta+\mf^2/\eta$ of the three-parton state at vanishing relative transverse momenta.

\subsection{Summary of notations}
\label{sec:not}

For the reader's convenience we collect here the notations that are used throughout. For the emission
by the antiquark,
\begin{align}
&\kt=\kk-\zeta\qq\,,\qquad \pt=\pp-\frac{\xi}{1-\zeta}(\qq-\kk)\,,\qquad
\ef^2=\zeta(1-\zeta)Q^2+\mf^2\,,\qquad \of=\frac{\xi\,\mf}{1-\zeta}\,,
\nonumber\\
&\Ef^2=Q^2+\mf^2\left(\frac1\zeta+\frac1\eta\right)\,,\qquad
a\equiv\frac{\xi\,\eta}{\zeta(1-\zeta)^2}\,.
\label{eq:notation}
\end{align}
The photon vertex and the vertex for the gluon emission by the antiquark are, respectively,
\begin{equation}
\varphi^{ij}_{\lambda_1\lambda}(\zeta)\,\kt^j+\mf\,\mu^i_{\lambda_1\lambda}\,,\qquad\qquad
\tau^{mn}_{\lambda\lambda_2}(\xi,\eta)\,\pt^{\,n}+\xi\,\of\,\mu^m_{\lambda\lambda_2}\,.
\label{eq:vertices}
\end{equation}
For the emission by the quark the roles of $\zeta$ and $\eta$ are exchanged, and we define
\begin{equation}
\eft^{\,2}\equiv\eta(1-\eta)Q^2+\mf^2\,,\qquad \oft\equiv\frac{\xi\,\mf}{1-\eta}\,,\qquad
\tilde a\equiv\frac{\xi\,\zeta}{\eta(1-\eta)^2}\,,
\label{eq:notation-q}
\end{equation}
while $\Ef$ is invariant under $\zeta\leftrightarrow\eta$.
